\documentclass[12pt]{article}
\usepackage[utf8]{inputenc}
\usepackage[T1]{fontenc}
\usepackage{hyperref}
\title{Sangu Language Resources for K-12 Curriculum Development}
\author{Compiled by Claude Sonnet 5.5 (claude-sonnet-5-5) for Andr\'as Kornai}
\date{}
\begin{document}
\maketitle

\section{Introduction}
Sangu (ISO 639-3: sbp; also Kisangu, Kisango, Kirori, Eshisango, Rori) is a Bantu language of Tanzania, spoken by the Sangu people of the Usangu Plain, Chunya District, Mbeya Region. A 1987 estimate gives about 75,000 speakers. Do not confuse it with Sangu of Gabon (snq), Sanga (sng) of the DRC, or Sango (sag). The search results show little published material: a Bible translation and audio recordings, wordlists, and one theological thesis. No primer, textbook or published dictionary was found.

\section{Primary Educational Resources}
No Sangu primers, textbooks or school programmes were found. The language is not known to be taught in schools.

\section{Linguistic and Reference Materials}
\begin{itemize}
\item Wikipedia, Sangu language (Tanzania): \url{https://en.wikipedia.org/wiki/Sangu_language_(Tanzania)}
\item Wikipedia, Sangu people: \url{https://en.wikipedia.org/wiki/Sangu_people}
\item ASJP wordlist for Sangu (sbp), based on Nurse and Philippson and the Tanzanian Language Survey: \url{https://asjp.clld.org/languages/SANGU}
\item SATS thesis on translating ``spirits'' in 1 Peter 3:19 into Sangu, with an inventory of candidate words: \url{https://research.sats.ac.za/items/c15ddddb-aa6a-41eb-972b-55349c2d0225}
\item SOAS Library manuscript record that appeared in the search: \url{https://library.soas.ac.uk/Record/MS_380565.01.02} (contents not verified)
\end{itemize}

\section{Religious and Community Resources}
Scripture Earth lists Sangu scripture resources (a Bible on YouVersion, a Luke film and an app): \url{https://www.scriptureearth.org/00eng.php?iso=sbp}

\section{Digital and Audio Resources}
\begin{itemize}
\item Global Recordings Network, Sangu: \url{https://globalrecordings.net/en/language/sbp}
\item Wikitongues language page (EGIDS 6b): \url{https://wikitongues.org/languages/sbp/}
\end{itemize}

\section{Curriculum Development Context}
Sangu has a sizeable speaker base but almost no published educational or lexical material found here. A curriculum would need new work in orthography, a dictionary and reading materials, ideally using the existing Bible translation as a text base and working with Sangu community members. Swahili is the language of schooling in Tanzania.
\end{document}